% Header: General study / work / project template
%%%%%%%%%%%%%%%%%%%%%%%%%%%%%%%%%%%%%%%%%%%%%%%%%%%%%%%%%%%%%%%%%%%

\documentclass[11pt,a4paper]{scrartcl}

% Layout
\usepackage[margin=2.5cm]{geometry}
\usepackage[onehalfspacing]{setspace}

% General packages
\usepackage{graphicx}
\usepackage{enumitem}
\usepackage{tcolorbox}
\tcbuselibrary{breakable}
\usepackage[breaklinks=true,colorlinks=true,linkcolor=blue,urlcolor=blue,citecolor=blue]{hyperref}
\usepackage{amsmath,amsthm,amssymb,mathtools}
\usepackage{braket}
\usepackage{icomma}
\usepackage[T1]{fontenc}
\usepackage[utf8]{inputenc}

% Language: choose one
% \usepackage[ngerman]{babel}
\usepackage[english]{babel}

% Units / tables / floats / references
\usepackage[locale = UK,space-before-unit=true,per-mode = symbol]{siunitx}
\usepackage{booktabs,multirow}
\usepackage{placeins}
\usepackage[natbib,abbreviate=true,doi=false,style=numeric-comp,giveninits=true,sorting=none]{biblatex}
\usepackage{csquotes}

% Optional bibliography
% \addbibresource{references.bib}

% Graphics paths
\graphicspath{{Bilder/} {images/} {figures/} {chapters/}}

% Optional math shortcuts
\newcommand{\R}{\mathbb{R}}
\newcommand{\N}{\mathbb{N}}
\newcommand{\Q}{\mathbb{Q}}
%\newcommand{\set}[1]{\left\{ #1 \right\}}
\newcommand{\abs}[1]{\left| #1 \right|}
\newcommand{\norm}[1]{\left\lVert #1 \right\rVert}
\newcommand{\ol}[1]{\overline{#1}}

% Optional units
\DeclareSIUnit{\dBm}{dBm}

% Box styles
\newtcolorbox{calculationbox}[1][]{
    title=Calculation,
    breakable,
    #1
}

\newtcolorbox{solutionbox}[1][]{
    title=Solution,
    breakable,
    #1
}

\newtcolorbox{answerbox}[1][]{
    title=Answer,
    breakable,
    #1
}

\newtcolorbox{notebox}[1][]{
    title=Notes,
    breakable,
    #1
}

\newtcolorbox{sidecalcbox}[1][]{
    title=Side Calculation,
    breakable,
    #1
}

%%%%%%%%%%%%%%%%%%%%%%%%%%%%%%%%%%%%%%%%%%%%%%%%%%%%%%%%%%%%%%%%%%%
\begin{document}

    \title{Theoretical Physics 2 - Blatt 06}
    \author{Tobias Laser}
    \date{\today}
    \maketitle
    \vfill
    \thispagestyle{empty}

    \pagenumbering{arabic}

    % ================================================================
    \paragraph{15. Orbital Angular Momentum Operatos and Spherical Harmonics} 

        We consider the orbital angular momentum operators $\hat{L}_j = \sum\limits_{k, l = 0}^3 \epsilon_{jkl} \hat{x}_k \hat{p}_l$.

        \begin{enumerate}
            \item 
                Write the operators $\hat{L}_j$, $\sum\limits_j\hat{L}^2_j$, $\hat{L}_\pm = \hat{L}_1 \pm i \hat{L}_2$ in the \textbf{in the position representation} (in cartesian coordinates).
                \begin{calculationbox}
                    We are given $\hat{L}_j = \sum\limits_{k, l = 1}^3 \epsilon_{j,k,l}\hat{x}_k \hat{p}_l$

                    In cartesian coordinates $\hat{p}_l = - i \hbar \frac{\partial}{\partial x_l}$

                    so we get $\hat{L}_j = - \hbar \sum\limits_{k,l = 1}^3 \epsilon_{jkl} \frac{\partial}{\partial x_L}$

                    Single component:

                    $j = 1$
                    \begin{calculationbox}
                        $\hat{L}_1 = \hat{x}_2 \hat{p}_3 - \hat{x}_3 \hat{p}_2 = - i \hbar (x_2 \frac{\partial}{\partial x_3} - x_3 \frac{\partial}{\partial_2})$
                    \end{calculationbox}

                    $j = 2$
                    \begin{calculationbox}
                        $\hat{L}_2 = \hat{x}_3 \hat{p}_1 - \hat{x}_1 \hat{p}_3 = - i \hbar (x_3 \frac{\partial}{\partial x_1} - x_1 \frac{\partial}{\partial_3})$
                    \end{calculationbox}

                    $j = 3$
                    \begin{calculationbox}
                        $\hat{L}_3 = \hat{x}_1 \hat{p}_2 - \hat{x}_2 \hat{p}_1 = - i \hbar (x_1 \frac{\partial}{\partial x_2} - x_2 \frac{\partial}{\partial_1})$
                    \end{calculationbox}

                    Total operator $\hat{L}^2$
                    \begin{calculationbox}
                        $\hat{L}^2 = \hat{L}_1^2 + \hat{L}_2^2 + \hat{L}_3^2 = -\hbar^2 
                        [ (x_2 \frac{\partial}{\partial x_3} - x_3 \frac{\partial}{\partial x_2})^2 
                        + (x_3 \frac{\partial}{\partial x_1} - x_1 \frac{\partial}{\partial x_3})^2 
                        + (x_1 \frac{\partial}{\partial x_2} - x_2 \frac{\partial}{\partial x_1})^2]$
                    \end{calculationbox}

                    Ladder operator 
                    \begin{calculationbox}
                        $\hat{L}_\pm = \hat{L}_2 \pm i \hat{L}_2 = 
                        - i \hbar (x_2 \frac{\partial}{\partial x_3} - x_3 \frac{\partial}{\partial x_2}) 
                        \pm \hbar (x_3 \frac{\partial}{\partial x_1} - x_1 \frac{\partial}{\partial x_3})$
                    \end{calculationbox}
                \end{calculationbox}
            \item 
                Express the same operators in \textbf{spherical coordinates}, using $x_1 = r \sin\theta \cos\phi$, $x_2 = r \sin\theta \sin\phi$, $x_3 = r \cos\theta$. (e.g., by applying the chain rule).

                \begin{calculationbox}
                    Chainrule $\frac{\partial}{\partial x_i} 
                    = \frac{\partial r}{\partial x_i} \frac{\partial}{\partial r} 
                    + \frac{\partial \theta}{ \partial x_i} \frac{\partial}{\partial \theta} 
                    + \frac{\partial \phi}{\partial x_i} \frac{\partial}{\partial \phi}$

                    Identities: $4 = \sqrt{x_1^2 + x_2^2 + x_3^2}$, $x_1^2 x_2^2 = r^2 \sin^2\theta$, $\phi = \arctan{\frac{x_2}{x_1}}$

                    derivations
                    \begin{calculationbox}
                        derivation of r
                        \begin{calculationbox}
                            $\frac{r}{x_1} = \frac{x_1}{3} = \sin\theta \cos \phi$
                        \end{calculationbox}

                        derivation of $\phi$
                        \begin{calculationbox}
                            $\frac{\partial \phi}{\partial x_1} 
                            = - \frac{x_2}{x_1^2 + x_2^2} 
                            = -\frac{r \sin\theta \cos\phi}{r^2 \sin^2\theta} 
                            = - \frac{\sin\phi}{r \sin\theta}$
                        \end{calculationbox}

                        derivation of $\theta$
                        \begin{calculationbox}
                            $\cos\theta = \frac{x_3}{r}$

                            $\frac{\partial}{\partial x_1} \cos\theta 
                            = -\sin\theta \frac{\partial}{\partial x_1} 
                            = -\frac{x_3}{r^2} \frac{\partial r}{\partial x_1}
                            = -\frac{x_3}{r^2} \frac{x_1}{r}
                            =  -\frac{r \cos\theta r \sin\theta \cos\phi}{r^3}
                            $

                            $\frac{\partial}{\partial x_1} = \frac{\cos\theta \cos \phi}{r}$
                        \end{calculationbox}
                    \end{calculationbox}

                    partial derivations by $x_i$
                    \begin{calculationbox}
                        $x_1$
                        \begin{calculationbox}
                            $\frac{\partial}{\partial x_1} 
                            = \sin\theta\cos\phi\phi \frac{\partial}{\partial r} 
                            + \frac{\cos \theta \cos \phi}{r} \frac{\partial}{\partial \theta}
                            +\frac{\cos\phi}{r \sin\theta} \frac{\partial}{\partial \phi}$
                        \end{calculationbox} 

                        $x_2$
                        \begin{calculationbox}
                            $\frac{\partial}{\partial x_2}
                            = \sin\theta\sin\phi \frac{\partial}{\partial r}
                            + \frac{\cos\theta\cos\phi}{r} \frac{\partial}{\partial \theta} 
                            - \frac{\cos\phi}{r \sin\theta} \frac{\partial}{\partial \phi}
                            $
                        \end{calculationbox}

                        $x_3$
                        \begin{calculationbox}
                            $\frac{\partial}{\partial x_3} 
                            = \cos\theta \frac{\partial}{\partial r} - \frac{\sin \theta}{ r} \frac{\partial}{\partial \theta}
                            $
                        \end{calculationbox}
                    \end{calculationbox}

                    Components 
                    \begin{calculationbox}
                        $\hat{L_1} = i\hbar
                        (\sin\phi \frac{\partial}{\partial \theta} 
                        + \cot\theta \cos \phi \frac{\partial}{\partial \phi})$

                        $\hat{L}_2 = i \hbar 
                        ( -\cos \phi \frac{\partial}{\partial \theta}
                        + \cot\theta \sin\phi \frac{\partial}{\partial \phi})$

                        $\hat{L}_3 = - i \hbar \frac{\partial}{\partial \phi}$
                    \end{calculationbox}

                    total operator $\hat{L}^2$
                    \begin{calculationbox}
                        $\hat{L}^2 = -\hbar [\frac{1}{\sin\theta} \frac{\partial}{\partial \theta} (\sin\theta \frac{\partial}{\partial \theta}) + \frac{1}{\sin^2 \theta} \frac{\partial^2}{\partial \phi^2}]$
                    \end{calculationbox}
                    \begin{calculationbox}
                        
                    \end{calculationbox}

                    Ladder operator
                    \begin{calculationbox}
                        $\hat{L}_\pm = \hbar e^{\pm i \phi} (\pm\frac{\partial}{\partial \theta} + i \cot\theta \frac{\partial}{\partial\phi})$
                    \end{calculationbox}
                    \begin{calculationbox}
                        
                    \end{calculationbox}
                \end{calculationbox}

            \item 
                Formulate the two differential equations that a common eigenstate $\psi_{k, l, m}(r, \theta, \phi)$ of
                \begin{itemize}
                    \item 
                        $\hat{L}^2$ (with eigenvalue $\hbar^2 l (l + 1)$)
                        \begin{calculationbox}
                            $\hat{L}^2 \psi_{k,l,m} 
                            = -\hbar^2 [\frac{1}{\sin\theta} \frac{\partial}{\partial \theta} (\sin\theta \frac{\partial \psi_{k,l,m}}{\partial \theta}) + \frac{1}{\sin^2 \theta} \frac{\partial^2 \psi_{k,l,m}}{\partial \phi^2}] 
                            = \hbar^2 l (l + 1) \psi_{k,l,m}
                            $

                            $-[\frac{1}{\sin\theta} \frac{\partial}{\partial \theta} (\sin\theta \frac{\partial \psi_{k,l,m}}{\partial \theta}) + \frac{1}{\sin^2 \theta} \frac{\partial^2 \psi_{k,l,m}}{\partial \phi^2}] 
                            = l (l + 1) \psi_{k,l,m}
                            $
                        \end{calculationbox}
                    \item 
                        $\hat{L}_3$ (with eigenvalue $\hbar m$)
                        \begin{calculationbox}
                            $\hat{L}_3 \psi_{k,l,m} 
                            = - i \hbar \frac{\partial \psi_{k,l,m}}{\partial \phi} 
                            = \hbar m \psi_{k,l,m}
                            $

                            $\frac{\partial \psi_{k,l,m}}{\partial \phi} = i m \psi_{k,l,m}$
                        \end{calculationbox}
                \end{itemize}
                must satisfy, \textbf{without solving them}

                Is it always possible to find a solution that satisfies both differential eqations simultaneously?
                \begin{calculationbox}
                    Yes, because commutator $[\hat{L}^2, \hat{L}_3] = 0$ and therefor have common eigenfunctions.
                    \begin{calculationbox}
                        $[\hat{L}^2, \hat{L}_3]$

                        $= [\hat{L}_1^2, \hat{L}_3] + [\hat{L}_2^2, \hat{L}_3] + [\hat{L}_3^2, \hat{L}_3]$

                        $= \hat{L}_1[\hat{L}_1, \hat{L}_3] + [\hat{L}_1, \hat{L}_3] \hat{L}_1 + \hat{L}_2[\hat{L}_2, \hat{L}_3] + [\hat{L}_2, \hat{L}_3] \hat{L}_2$

                        $= \hat{L}_1(-i\hbar \hat{L}_2) + (-i \hbar \hat{L}_2) \hat{L}_1 + \hat{L}_2(i\hbar \hat{L}_1) + (i\hbar \hat{L}_1) \hat{L}_2$
                        
                        $= 0$
                    \end{calculationbox}
                \end{calculationbox}
            
            \item 
                The eigenstate $\psi_{k, l, l}(r, \theta, \psi)$ satisfies $\hat{L}_+ \psi_{k, l, l} = 0$.
                \begin{itemize}
                    \item 
                        Consider the resulting differential equation for the special case $l = 1$, and sovle it.
                        \begin{calculationbox}
                            We use $\hat{L}_+$ and $\hat{L}_-$ from (2)

                            ($_+ Y_l^m = \hbar \sqrt{l(l + 1) - m (m + 1)} Y_l^m + 1$)

                            For $l = 1$ holds true: $\hat{L}_+ Y_1^1 = 0$ and $Y_1^1$ is eigenfunction from $_3$ with $m = 1$: $Y_1^1(\theta, \phi) = f(\theta) \exp(i \phi)$

                            so we plug it in: 

                            $0 = \hat{L}_+ Y_1^1$

                            $= \hbar \exp(i \phi) (\frac{\partial}{\partial \theta} + i \cot\theta \frac{\partial}{\partial \phi})f(\theta)\exp(i\phi)$

                            $= \hbar \exp(i \phi)(f'(\theta) \exp(i\phi) + i \cot\theta i f(\theta) \exp{i\phi})$

                            So we get: $f'(\theta) - \cot\theta f(\theta) = 0$

                            $\frac{f'(\theta)}{f(\theta)} = \cot\theta$

                            Integrate: $\ln{f(\theta)} = \ln(\sin\theta) + C$

                            $f(\theta) = C \sin\theta$

                            Therefore $Y_1^1(\theta, \phi) = C \sin\theta \exp(i\phi)$
                        \end{calculationbox}

                    \item
                        The, using the operator $\hat{L}_-$, derive an explicit form of the spherical harmonics $Y_l^m(\theta, \phi)$, which are defined via $\psi_{k, l, m}(r, \theta, \phi) = R_k(r) Y_l^m(\theta, \phi)$

                        Here $R_k(r)$ remeins unspecified but is normalized via $\int\limits_0^\infty r^2 dr \abs{R_(r)}^2 = 1$.

                        \begin{calculationbox}
                            $\hat{L}_- Y_l^m = \hbar \sqrt{l(l+1) - m (m - 1)} Y_l^{m-1}$

                            for $l=1, m=1$

                            \begin{calculationbox}
                                $\hbar \sqrt{2} Y_1^0 = \hat{L}_-Y_1^1 $

                                $= \hbar \exp(-i \phi) (-\frac{\partial}{\partial \theta} + i \cot\theta \frac{\partial }{\partial \phi}) C \sin\theta \exp(i \phi)$

                                $= \hbar \exp(- i \phi) (- C \cos\theta \exp(i\phi) - C \cos\theta \exp(i\phi))$

                                $= - 2\hbar C \cos\theta$

                                Therefore $y_1^0 = -\sqrt{2} C \cos\theta$ with $C = -\sqrt{\frac{3}{8\pi}}$
                            \end{calculationbox}

                            for $l=1, m=-1$
                            \begin{calculationbox}
                                $\hbar \sqrt{2} Y_1^{-2} = \hat{L}_- Y_1^0$

                                $= \hbar \exp(-i \phi) (-\frac{\partial}{\partial \theta} + i \cot\theta \frac{\partial }{\partial \phi}) \sqrt{\frac{3}{4\pi}} \cos\theta$

                                $= \hbar \exp(- i \phi) (- \sqrt{\frac{3}{4\pi}} \sin\theta)$

                                Therefore $Y_1^{-1} = \sqrt{\frac{3}{8\pi}} \sin\theta \exp{-i \phi}$
                            \end{calculationbox}
                        \end{calculationbox}
                \end{itemize}
            
            \item
                Insert the spherical harmonics into the differential equations from part (3), and show that they are indeed the corresponding eigenfunctions.
                \begin{calculationbox}
                    Previously $\hat{L}^2 Y_1^- = \hbar^2 1 (1+ 1) Y_1^m = 2 \hbar^2 Y_1^m$ and $\hat{L}_3Y_1^m = \hbar m Y_1^m$

                    Check with $\hat{L}_3$
                    \begin{calculationbox}
                        For $\hat{L}_3 = - i \hbar \frac{\partial}{\partial \phi}$ for $Y_1^1 = -\sqrt{\frac{3}{8\pi}}\sin\theta \exp{i\phi}$
                        
                        We get $\frac{\partial}{\partial \phi} Y_1^1 = i Y_1^1$

                        therefor $\hat{L}_3 Y_1^1 = - i \hbar i Y_1^1 = \hbar Y_1^1$ So $m = 1$

                    \end{calculationbox}

                    \begin{calculationbox}
                        For $Y_1^0 = \sqrt{\frac{3}{4\pi}} \cos\theta$ 
                        
                        We get $\frac{\partial}{\partial \phi} Y_1^0 = 0$

                        therefor $\hat{L}_3 Y_1^0 = 0$ So $m = 0$

                    \end{calculationbox}

                    \begin{calculationbox}
                        $Y_1^{-1} = \sqrt{\frac{3}{8\pi}} \sin\theta \exp{-i\phi}$
                        
                        We get $\frac{\partial}{\partial \phi} Y_1^{-1} = -i Y_1^{-1}$

                        therefor $\hat{L}_3 Y_1^{-1} = -i \hbar (-i) Y_1^{-1} = -\hbar Y_1^{-1}$ So $m = -1$

                    \end{calculationbox}

                    Check with $\hat{L}^2$
                    \begin{calculationbox}

                        From (2) $\hat{L}^2 = -\hbar^2 [\frac{1}{\sin\theta}\frac{\partial}{\partial \theta}(\sin\theta \frac{\partial}{\partial \theta}) + \frac{1}{\sin^2 \theta}\frac{\partial^2}{\partial\phi^2}]$

                        for $Y_1^0 = A \cos\theta$, $A =\sqrt{\frac{3}{4\pi}}$

                        $\frac{\partial^2}{\partial \phi^2} Y_1^0 = 0$

                        $\frac{\partial}{\partial \theta} Y_1^0 = -A \sin\theta$
                        
                        $\sin\theta \frac{\partial}{\partial \theta} Y_1^0 = -A \sin^2\theta$

                        $\frac{\partial}{\partial \theta} (\sin\theta \frac{\partial}{\partial \theta} Y_1^0) = -2 A \sin\theta \cos\theta$

                        $\frac{1}{\sin\theta}\frac{\partial}{\partial \theta} (\sin\theta \frac{\partial}{\partial \theta} Y_1^0) = -2 A \cos\theta = -2 Y_1^0$

                        So $\hat{L}^2 Y_1^0 = -\hbar^2 (-2 Y_1^0) = 2 \hbar^2 Y_1^0$

                        For $Y_1^{\pm1}$ we get the same expression $\hat{L}^2 Y_1^{\pm1} = s\hbar^2 Y_1^{\pm1}$

                        Therefore $\hat{L} Y_1^m = 2 \hbar^2 Y_1^m$ and $l (l+1) = 2$ so $l = 1$
                    \end{calculationbox}

                \end{calculationbox}

        \end{enumerate}

    \paragraph{16. Scattering from a Box Ootential, Bound States, and Scattering Resonances} We consider a three-dimensional Potential $V(r)$ (in spherical coordinates) 
    $V(r) = \begin{cases}
        -V_0 & r \leq a\\ 0 & r > a
    \end{cases}$
    with a positive constant $V_0$.

    \begin{enumerate}
        \item 
            Formulate the \textbf{radial Schrödinger equation} for the two regions $r < a$ and $r > a$ in the case $l = 0$ (see Exercise 15), and find the \textbf{general solution for} $E > 0$ using the condition $P_l(r \to 0) = 0$, where $P_l$ is the radial part of the wavefunction, i.e. $\psi_{E,l = 0, m = 0}(\vec{x}) = \frac{P_{l = 0}(r)}{r} >_0^0(\theta, \phi)$. Furthermore, using the condition $\left.\frac{1}{P_0(r)} \frac{d P_0(r)}{dr}\right|_{r = a-0} = \left. \frac{1}{P_0(r)} \frac{d P_0 (r)}{dr} \right|_{r=a+0}$, derive an equation for $\tan\delta_0$, where $\delta_0$ is the $scattering phase shift for$ $l=0$.

            \begin{calculationbox}

                For $\ell = 0$ we use
                \[
                \psi_{E,0,0}(\mathbf{x}) = \frac{P_0(r)}{r} Y_0^0(\theta,\phi)
                \]

                The radial Schrödinger equation becomes
                \begin{calculationbox}
                    $-\frac{\hbar^2}{2m}\frac{d^2 P_0}{dr^2} + V(r)P_0 = E P_0$
                \end{calculationbox}

                Region I: $r < a$
                \begin{calculationbox}
                    $V(r) = -V_0$

                    $-\frac{\hbar^2}{2m}P_0'' - V_0 P_0 = E P_0$

                    $P_0'' + \frac{2m(E+V_0)}{\hbar^2} P_0 = 0$

                    Define $K^2 = \frac{2m(E+V_0)}{\hbar^2}$

                    $P_0'' + K^2 P_0 = 0$

                    General solution:

                    $P_0^{<}(r) = A \sin(Kr) + B \cos(Kr)$

                    Using $P_0(r \to 0)=0$:

                    $B = 0$

                    $\Rightarrow P_0^{<}(r) = A \sin(Kr)$
                \end{calculationbox}

                Region II: $r > a$
                \begin{calculationbox}
                    $V(r) = 0$

                    $P_0'' + \frac{2mE}{\hbar^2} P_0 = 0$

                    Define $k^2 = \frac{2mE}{\hbar^2}$

                    $P_0'' + k^2 P_0 = 0$

                    Scattering solution:

                    $P_0^{>}(r) = B \sin(kr + \delta_0)$
                \end{calculationbox}

                Matching condition at $r = a$
                \begin{calculationbox}
                    $\left.\frac{P_0'}{P_0}\right|_{a-0}
                    =
                    \left.\frac{P_0'}{P_0}\right|_{a+0}$
                \end{calculationbox}

                Inside:
                \begin{calculationbox}
                    $P_0^{<} = A \sin(Kr)$

                    $P_0' = AK \cos(Kr)$

                    $\frac{P_0'}{P_0} = K \cot(Kr)$

                    $\Rightarrow \left.\frac{P_0'}{P_0}\right|_{a-0} = K \cot(Ka)$
                \end{calculationbox}

                Outside:
                \begin{calculationbox}
                    $P_0^{>} = B \sin(kr + \delta_0)$

                    $P_0' = Bk \cos(kr + \delta_0)$

                    $\frac{P_0'}{P_0} = k \cot(kr + \delta_0)$

                    $\Rightarrow \left.\frac{P_0'}{P_0}\right|_{a+0} = k \cot(ka + \delta_0)$
                \end{calculationbox}

                Matching gives:
                \begin{calculationbox}
                    $K \cot(Ka) = k \cot(ka + \delta_0)$
                \end{calculationbox}

                Solve for $\tan\delta_0$:
                \begin{calculationbox}
                    $\tan(ka + \delta_0) = \frac{k}{K}\tan(Ka)$

                    Using
                    $\tan(ka + \delta_0)
                    =
                    \frac{\tan(ka) + \tan\delta_0}{1 - \tan(ka)\tan\delta_0}$

                    we obtain

                    $\tan\delta_0
                    =
                    \frac{\frac{k}{K}\tan(Ka) - \tan(ka)}
                    {1 + \frac{k}{K}\tan(Ka)\tan(ka)}$
                \end{calculationbox}

                with
                \begin{calculationbox}
                    $k = \sqrt{\frac{2mE}{\hbar^2}}, \quad
                    K = \sqrt{\frac{2m(E+V_0)}{\hbar^2}}$
                \end{calculationbox}

            \end{calculationbox}

        \item 
            \textbf{Scattering at high and low energies}
            \begin{enumerate}
                \item
                    What is $\delta_0$ in the limit $E \to \infty$? What is the physical meaning?

                    \begin{calculationbox}
                        From part (a) we have
                        \[
                        \tan\delta_0
                        =
                        \frac{\frac{k}{K}\tan(Ka)-\tan(ka)}
                        {1+\frac{k}{K}\tan(Ka)\tan(ka)}
                        \]

                        with
                        \[
                        k = \sqrt{\frac{2mE}{\hbar^2}},
                        \qquad
                        K = \sqrt{\frac{2m(E+V_0)}{\hbar^2}}.
                        \]

                        In the limit $E \to \infty$, the potential depth $V_0$ becomes negligible compared to $E$:
                        \[
                        E+V_0 \approx E.
                        \]

                        Therefore
                        \[
                        K \approx k.
                        \]

                        Hence
                        \[
                        \frac{k}{K} \to 1,
                        \qquad
                        Ka \to ka.
                        \]

                        Insert this into the expression for $\tan\delta_0$:
                        \[
                        \tan\delta_0
                        =
                        \frac{\tan(ka)-\tan(ka)}
                        {1+\tan^2(ka)}
                        =
                        0.
                        \]

                        Thus
                        \[
                        \boxed{\delta_0 \to 0}
                        \]
                        for $E \to \infty$.
                    \end{calculationbox}

                    Physical interpretation:
                    \begin{calculationbox}
                        At very high energy, the particle is barely affected by the finite potential well.

                        The wavelength is very small compared to the length scale of the potential, and the scattering phase shift vanishes.

                        Therefore, the potential becomes effectively transparent:
                        \[
                        \boxed{\delta_0 \to 0}
                        \]
                        means almost no scattering.
                    \end{calculationbox}

                \item 
                     Determine $\delta_0$ in the limit $E \to 0$, and calculate the \textbf{energy-dependent total cross section} $\sigma(E)$

                     Find the specific values of $V_0$ and $a$ that lead to \textbf{scattering resonances}.

                     Show that for small energies, the scattering process reduces to \textbf{s-wave scattering} (i.e., $l=0$).
                     \begin{calculationbox}
                        From part (1):
                        \[
                        \tan\delta_0
                        =
                        \frac{\frac{k}{K}\tan(Ka)-\tan(ka)}
                        {1+\frac{k}{K}\tan(Ka)\tan(ka)}
                        \]

                        with
                        \[
                        k=\sqrt{\frac{2mE}{\hbar^2}},
                        \qquad
                        K=\sqrt{\frac{2m(E+V_0)}{\hbar^2}}.
                        \]

                        For $E \to 0$ we have
                        \[
                        k \to 0,
                        \qquad
                        K \to \kappa,
                        \qquad
                        \kappa := \sqrt{\frac{2mV_0}{\hbar^2}}.
                        \]

                        Also
                        \[
                        \tan(ka) \approx ka.
                        \]

                        Therefore
                        \[
                        \tan\delta_0
                        \approx
                        \frac{\frac{k}{\kappa}\tan(\kappa a)-ka}
                        {1}
                        \]

                        so
                        \[
                        \tan\delta_0
                        \approx
                        k\left(
                        \frac{\tan(\kappa a)}{\kappa}-a
                        \right).
                        \]

                        Define the scattering length
                        \[
                        a_s
                        :=
                        a-\frac{\tan(\kappa a)}{\kappa}.
                        \]

                        Then
                        \[
                        \boxed{
                        \tan\delta_0 \approx -ka_s
                        }
                        \]

                        and for small $k$:
                        \[
                        \boxed{
                        \delta_0 \approx -ka_s
                        }.
                        \]
                    \end{calculationbox}

                    Total cross section:
                    \begin{calculationbox}
                        For s-wave scattering,
                        \[
                        \sigma(E)
                        =
                        \frac{4\pi}{k^2}\sin^2\delta_0.
                        \]

                        Using
                        \[
                        \tan\delta_0 \approx -ka_s
                        \]
                        and
                        \[
                        \sin^2\delta_0
                        =
                        \frac{\tan^2\delta_0}{1+\tan^2\delta_0},
                        \]
                        we get
                        \[
                        \sigma(E)
                        =
                        \frac{4\pi}{k^2}
                        \frac{k^2a_s^2}{1+k^2a_s^2}.
                        \]

                        Therefore
                        \[
                        \boxed{
                        \sigma(E)
                        =
                        \frac{4\pi a_s^2}{1+k^2a_s^2}
                        }
                        \]

                        with
                        \[
                        a_s
                        =
                        a-\frac{\tan(\kappa a)}{\kappa},
                        \qquad
                        \kappa=\sqrt{\frac{2mV_0}{\hbar^2}}.
                        \]
                    \end{calculationbox}

                    Scattering resonances:
                    \begin{calculationbox}
                        A scattering resonance occurs when the scattering length diverges:
                        \[
                        a_s \to \infty.
                        \]

                        This happens when
                        \[
                        \tan(\kappa a) \to \infty.
                        \]

                        Therefore
                        \[
                        \kappa a
                        =
                        \left(n+\frac{1}{2}\right)\pi,
                        \qquad
                        n=0,1,2,\dots
                        \]

                        Since
                        \[
                        \kappa^2=\frac{2mV_0}{\hbar^2},
                        \]
                        this gives
                        \[
                        \boxed{
                        V_0a^2
                        =
                        \frac{\hbar^2}{2m}
                        \left(n+\frac{1}{2}\right)^2\pi^2
                        }
                        \]

                        At resonance,
                        \[
                        \sigma(E)
                        \to
                        \frac{4\pi}{k^2}.
                        \]
                    \end{calculationbox}

                    Low-energy scattering is s-wave scattering:
                    \begin{calculationbox}
                        For small energies, higher partial waves are suppressed.

                        In general,
                        \[
                        \delta_l \sim k^{2l+1}.
                        \]

                        Therefore,
                        for $l>0$,
                        \[
                        \delta_l \to 0
                        \quad
                        \text{faster than}
                        \quad
                        \delta_0.
                        \]

                        Hence only the $l=0$ contribution remains relevant:
                        \[
                        \boxed{
                        \text{low-energy scattering reduces to s-wave scattering.}
                        }
                        \]
                    \end{calculationbox}

                \item
                    Show that for certain values of $F_0 a^2$, the box potential becomes \textbf{transparent}, i.e. $\sigma = 0$ (\textbf{Ramsauer-Townsed effect}).
                    \begin{calculationbox}
                        From part (b)(ii), for low energies we found
                        \[
                        \sigma(E)
                        =
                        \frac{4\pi a_s^2}{1+k^2a_s^2}
                        \]
                        with
                        \[
                        a_s
                        =
                        a-\frac{\tan(\kappa a)}{\kappa},
                        \qquad
                        \kappa=\sqrt{\frac{2mV_0}{\hbar^2}}.
                        \]

                        The potential is transparent if
                        \[
                        \sigma(E)=0.
                        \]

                        This happens if
                        \[
                        a_s=0.
                        \]

                        Therefore
                        \[
                        a-\frac{\tan(\kappa a)}{\kappa}=0.
                        \]

                        Hence
                        \[
                        \tan(\kappa a)=\kappa a.
                        \]

                        Thus the transparency condition is
                        \[
                        \boxed{
                        \tan(\kappa a)=\kappa a
                        }
                        \]

                        with
                        \[
                        \kappa a
                        =
                        a\sqrt{\frac{2mV_0}{\hbar^2}}.
                        \]

                        Therefore the corresponding values of $V_0a^2$ are determined by
                        \[
                        \boxed{
                        V_0a^2
                        =
                        \frac{\hbar^2}{2m}x_n^2
                        }
                        \]
                        where $x_n$ are the non-zero solutions of
                        \[
                        \tan x_n = x_n.
                        \]
                    \end{calculationbox}

                    Physical interpretation:
                    \begin{calculationbox}
                        For these special values of $V_0a^2$, the scattered wave destructively interferes.

                        Therefore the phase shift effectively gives no scattering contribution, and the total cross section vanishes:
                        \[
                        \boxed{
                        \sigma = 0
                        }
                        \]

                        This is the Ramsauer--Townsend effect.
                    \end{calculationbox}
            \end{enumerate}
        
        \item 
            \textbf{Bound states for s-wave scattering}
            \begin{enumerate}
                \item 
                    Show  that the general solution of the radial Schrödinger equation for $E < 0$ ha the form $T_{t=0}(r) = \begin{cases}
                        c_1 \sin(Kr) & r \leq a\\ c_2 \exp(-\rho r) & r > a
                    \end{cases}$
                    and determine $\rho$ and $K$.
                \item 
                    Using the continuity conditions at $r = a$, show that $\rho$ satisfies $\tan(K a) = \frac{K}{\rho}$.
                    Discuss this equation:
                    \begin{itemize}
                        \item 
                            Determine the number of bound s-states as a function of the potential depht (for fixed $a$)
                        \item 
                            Show in particula that no bound states exist if the potential is too shallow
                        \item
                            Finally, compare the bound states for $E \to 0^-$ with the scattering resonances in the case $E \to 0^+$ from part (2)(b)
                    \end{itemize}
            \end{enumerate}
    \end{enumerate}

\end{document}